\documentclass[11pt]{article}

\usepackage[T1]{fontenc}
\usepackage[utf8]{inputenc}
\usepackage[a4paper,margin=1in]{geometry}
\usepackage{amsmath,amssymb}
\usepackage{booktabs,longtable,array}
\usepackage{enumitem}
\usepackage{xcolor}
\usepackage[hidelinks]{hyperref}

\definecolor{codeblue}{RGB}{25,70,120}
\newcommand{\code}[1]{\textcolor{codeblue}{\texttt{#1}}}
\newcommand{\R}{\mathbb{R}}
\newcommand{\N}{\mathbb{N}}
\newcommand{\eps}{\varepsilon}
\newcommand{\argminlex}{\mathop{\mathrm{lex\mbox{-}arg\,min}}}

\title{Current MLEvolve Scheduler Decision Logic\\
       \large A Mathematical and Step-by-Step Description}
\author{MLEvolve Engineering Documentation}
\date{Implementation revision \code{2d845eb} --- 29 July 2026}

\begin{document}

\maketitle

\begin{abstract}
This report formalizes the decision logic implemented by the local single-GPU
scheduler.  The current example configuration selects
\code{parallel\_time\_aware}, whose central decision is a rolling-horizon
minimization of normalized, priority-weighted predicted flow time.  Before that
objective is evaluated, the scheduler applies hard constraints for queue
eligibility, starvation, exclusive probes, concurrency, backend support,
historical pair compatibility, batch-indexed runtime evidence, predicted GPU
memory, and live-memory admission.  This document describes those operations in
execution order, gives the exact objective and deterministic tie-break key, and
then summarizes the alternative legacy mode branches.
\end{abstract}

\tableofcontents

\section{Scope and implementation status}

The scheduler is a non-preemptive, rolling-horizon dispatcher for one physical
GPU.  On every service-loop iteration it first processes worker exits and user
commands, updates cache and GPU telemetry, and then attempts to dispatch pending
work.  In the current \code{parallel\_time\_aware} mode, an active training job
is never paused merely because a higher-priority job or an exclusive probe has
arrived.  New work is considered at the next legal decision boundary, while an
exclusive probe waits for the current work to drain.

The implemented objective version is
\code{time\_v3\_flow\_only}.  Despite the field name
\code{objective\_score}, a smaller value is better in this mode.  GPU/SM
utilization, predicted co-run slowdown, makespan rewards, throughput rewards,
and cached objectives from older policy versions do not affect this selection.

The main implementation is split across the following files:

\begin{itemize}[leftmargin=*]
  \item \code{scheduler/service.py}: the polling, drain, admission, dispatch,
        and runtime-fallback control flow;
  \item \code{scheduler/placement\_planner.py}: queue-window construction,
        mode branching, anchor selection, and final plan construction;
  \item \code{scheduler/candidate\_generator.py}: group, backend, and batch
        candidates;
  \item \code{scheduler/resource\_estimator.py}: batch-indexed average-VRAM and
        runtime estimates plus Pareto pruning;
  \item \code{scheduler/time\_objective.py}: feasibility checks, rolling-horizon
        flow cost, normalized score, search, and tie-breaking;
  \item \code{scheduler/compatibility.py}: packing eligibility and historical
        pair rejection;
  \item \code{scheduler/telemetry.py}: live-memory hysteresis; and
  \item \code{scheduler/supervisor.py}: the final backend and overlap gate.
\end{itemize}

\section{Notation}

At a scheduling instant $t$, let:

\begin{longtable}{>{$}l<{$} p{0.74\textwidth}}
\toprule
\text{Symbol} & Meaning \\
\midrule
\endhead
Q(t) & Set of runnable jobs, excluding jobs that are already active. \\
A(t) & Set of jobs currently running on the GPU. \\
W(t) & Bounded planning window formed from priority, age, starvation, and probe subsets. \\
G & A candidate group of new jobs, with $G\subseteq W(t)$. \\
h & An execution backend, for example \code{cuda\_process}, \code{mps}, \code{stream}, or \code{exclusive}. \\
p_i & User-supplied base priority of job $i$; larger values are more urgent. \\
q_i & Durable FIFO queue sequence of job $i$; smaller values are older in FIFO order. \\
s_i & Submission timestamp of job $i$. \\
b_i^0 & Immutable originally requested batch size of job $i$. \\
b_i & A candidate or selected batch size. \\
\widehat m_{i,h}(b_i) & Predicted average GPU memory in MiB for job $i$, backend $h$, and batch $b_i$. \\
\widehat e_{i,h}(b_i) & Predicted seconds per epoch. \\
r_i & Remaining epoch count. \\
\widehat p_{i,h}(b_i) & Predicted remaining runtime, equal to $r_i\widehat e_{i,h}(b_i)$. \\
B_{\mathrm{pred}} & Predicted-memory admission budget in MiB. \\
\eta & Priority-weight coefficient, currently $0.10$ in the example configuration. \\
C & Optional cap on the total number of parallel jobs; $C=\infty$ when unset. \\
\tau_{\mathrm{starve}} & Starvation timeout, currently $1800$ seconds in the example configuration. \\
\bottomrule
\end{longtable}

All comparisons described below are deterministic once queue state, profiles,
backend availability, telemetry state, and time $t$ are fixed.

\section{Top-level decision cycle}

One scheduler iteration performs the following ordered operations:

\begin{enumerate}[label=\textbf{Step \arabic*:},leftmargin=*]
  \item Poll active workers and record completed, failed, paused, or cancelled
        exits.  A packed group that loses members is reduced to its surviving
        members; one survivor is relabelled as an exclusive-shaped run.
  \item Consume pending submit, pause, resume, cancel, and preload commands.
        Only jobs in runnable persistent states enter $Q(t)$; recoverable jobs
        are omitted unless automatic resume is enabled.
  \item Warm likely baseline models.  This changes cache state but not the
        mathematical placement score.
  \item If work is active, sample device telemetry and update the live-memory
        admission gate described in Section~\ref{sec:live-gate}.
  \item Apply preemption logic.  This is deliberately a no-op in
        \code{parallel\_time\_aware}; older modes may pause a single exclusive
        job for a strictly higher effective-priority candidate.
  \item Reserve an exclusive probe if necessary, or construct a placement plan.
  \item Ask the worker supervisor to validate backend availability and overlap
        legality, apply selected batch overrides, launch the workers, and persist
        the full objective breakdown.
  \item In a concurrent mode, repeat the planning step immediately while another
        non-exclusive plan can legally be admitted.  Stop after an exclusive
        dispatch, a failed dispatch, or a no-plan result.
\end{enumerate}

This ordering matters.  For example, worker exits free capacity before the next
plan is scored, and the newest telemetry sample may close new admission without
interrupting already-running work.

\section{Time-aware placement algorithm}

\subsection{Step 1: effective priority and queue ordering}

For time-aware window construction, the effective priority is

\begin{equation}
  \pi_i(t)
  = p_i
  + \mathbf{1}_{\mathrm{aging}}
    \left\lfloor
      \frac{\max(0,t-s_i)}{T_{\mathrm{age}}}
    \right\rfloor \Delta_{\mathrm{age}},
  \label{eq:effective-priority}
\end{equation}

where $T_{\mathrm{age}}$ is the aging interval and
$\Delta_{\mathrm{age}}$ is the configured priority increment.  If aging is
disabled, if the interval is non-positive, or if the submission timestamp is
unavailable, the added term is zero.  The priority ordering key is

\begin{equation}
  k_i^{\mathrm{priority}}(t)=\bigl(-\pi_i(t),q_i,\mathrm{id}_i\bigr),
  \label{eq:priority-key}
\end{equation}

so higher effective priority wins, followed by FIFO sequence and job identifier.

The generic queue policy also uses priority-first FIFO ordering.  The
time-aware planner recomputes Equation~\eqref{eq:effective-priority} directly
from the submission timestamp when building its bounded window.

\subsection{Step 2: construct the rolling planning window}

Let $Q_N(t)$ contain normal jobs.  If the exclusive-probe feature is disabled,
jobs labelled as probes are treated as normal for this purpose.  Construct:

\begin{align}
  P_K(t) &= \text{first }K\text{ jobs of }Q_N(t)
            \text{ under }k^{\mathrm{priority}}, \\
  O_L(t) &= \text{oldest }L\text{ jobs under }(s_i,q_i,\mathrm{id}_i), \\
  S(t)   &= \{i\in Q_N(t):t-s_i\ge \tau_{\mathrm{starve}}\}, \\
  X(t)   &= \{i\in Q(t):i\text{ is an enabled exclusive probe}\}.
\end{align}

The deduplicated planning window is

\begin{equation}
  W(t)=\operatorname{sort}_{k^{\mathrm{priority}}}
  \bigl(P_K(t)\cup O_L(t)\cup S(t)\cup X(t)\bigr).
  \label{eq:window}
\end{equation}

The current example values are $K=8$ and $L=4$.  Notice that all starving jobs
and all enabled probes are included even if this makes $|W(t)|>K+L$.

If $S(t)\ne\varnothing$, define the mandatory anchor

\begin{equation}
  a^*(t)=\arg\min_{i\in S(t)}(s_i,q_i,\mathrm{id}_i).
  \label{eq:mandatory-anchor}
\end{equation}

Every candidate group must contain $a^*(t)$.  Otherwise, the ordinary anchor is
the first job in $W(t)$.  This rule is the scheduler's hard starvation
protection: urgency can change objective weights, but starvation changes the
feasible set itself.

\subsection{Step 3: exclusive-probe drain rule}

An enabled job with scheduling class \code{exclusive\_probe} reserves the next
idle boundary.  The service chooses one runnable probe by priority-FIFO order
and records its identifier durably.  Then:

\begin{itemize}[leftmargin=*]
  \item if $A(t)\ne\varnothing$, the planner returns no new dispatch and active
        work drains without preemption;
  \item once $A(t)=\varnothing$, the runnable set is restricted to the reserved
        probe, which is dispatched alone on \code{exclusive}; and
  \item normal admission resumes only after the reservation is cleared when the
        probe finishes or disappears from the runnable queue.
\end{itemize}

Thus a probe is a control-plane barrier, not merely a high objective weight.

\subsection{Step 4: concurrency and live-admission gates}

For an ordinary time-aware decision, let the remaining job slots be

\begin{equation}
  c_{\mathrm{new}}=
  \begin{cases}
    \infty, & C=\infty,\\
    C-|A(t)|, & C<\infty.
  \end{cases}
  \label{eq:remaining-slots}
\end{equation}

No candidate is generated when $c_{\mathrm{new}}\le 0$.  If active jobs exist,
new admission is also rejected when either the live-memory gate is closed or an
exclusive drain has been requested.  Precisely, the blocking condition is

\begin{equation}
  |A(t)|>0 \;\land\;
  \bigl(\neg\mathrm{AdmissionOpen}(t)
        \;\lor\;\mathrm{ExclusiveDrain}(t)\bigr).
  \label{eq:active-block}
\end{equation}

The implementation applies the live gate only to adding work alongside an
active run.  A closed gate does not, by itself, cancel active jobs or prevent an
idle GPU from making its first dispatch.

\subsection{Step 5: generate candidate job groups}

Let $u=\min(|W(t)|,c_{\mathrm{new}})$.  When
$u\le N_{\mathrm{exact}}$ (currently $N_{\mathrm{exact}}=3$), the candidate
family is the exact non-empty subset family

\begin{equation}
  \mathcal G={G\subseteq W(t):1\le |G|\le u\}.
  \label{eq:exact-groups}
\end{equation}

If there is a mandatory anchor, retain only groups containing it.

For larger windows, the implementation uses a deterministic beam.  Starting
from the empty state, or from $(a^*)$ when starvation supplies an anchor, each
job in window order is either omitted or appended while $|G|<u$.  Duplicate
states are removed, and only $B_{\mathrm{beam}}$ states are retained by the key

\begin{equation}
  k^{\mathrm{group\mbox{-}beam}}(G)
  =\bigl(-|G|,(\mathrm{id}_i)_{i\in G}\bigr),
  \label{eq:group-beam}
\end{equation}

with $B_{\mathrm{beam}}=64$ in the example configuration.  All retained
non-empty intermediate states are emitted.  The beam is therefore bounded and
deterministic, but it is an approximation rather than an exhaustive optimizer.

\subsection{Step 6: enumerate legal backends}

If the GPU is idle and $|G|=1$, the planner evaluates the job on
\code{exclusive}.  Otherwise it considers each non-exclusive backend $h$ in
configured priority order for which:

\begin{equation}
  \mathrm{Available}(h)=1
  \quad\land\quad
  \forall i\in G:\;\mathrm{AllowsBackend}(i,h)=1.
  \label{eq:backend-filter}
\end{equation}

The supervisor later repeats the physical availability check.  When overlapping
an active group, both the new and active backends must be non-exclusive and
listed as overlap-capable for time-aware scheduling.  Exclusive groups never
overlap any active group.

\subsection{Step 7: generate and prune batch options}

Time-aware scheduling uses the immutable requested batch $b_i^0$, not a batch
override selected by an earlier dispatch.  It requires $b_i^0=2^{k_i}$ for an
integer $k_i\ge0$.  With the configured exponent offsets
$\Delta=\{-2,-1,0,1,2\}$, the pre-clipping proposals are

\begin{equation}
  \widetilde{\mathcal B}_i
  =\left[2^{\max(0,k_i+\delta)}:\delta\in\Delta\right].
  \label{eq:batch-proposals}
\end{equation}

Each proposal is clipped to an optional job/global maximum $b_i^{\max}$
(take $b_i^{\max}=\infty$ when it is unset),
clamped below by one, and deduplicated without changing order:

\begin{equation}
  \mathcal B_i
  =\operatorname{unique}
    \left(\max\{1,\min(\widetilde b,b_i^{\max})\}:
    \widetilde b\in\widetilde{\mathcal B}_i\right).
  \label{eq:batch-clipping}
\end{equation}

For each $b\in\mathcal B_i$, the estimator must obtain both positive average
memory and positive epoch time.  With total epochs $E_i$ and completed epochs
$e_i^{\mathrm{done}}$, it computes

\begin{align}
  r_i &= \max(0,E_i-e_i^{\mathrm{done}}), \\
  \widehat p_{i,h}(b) &= r_i\widehat e_{i,h}(b).
  \label{eq:remaining-runtime}
\end{align}

An option is discarded if the total epoch count, average memory, or runtime is
missing.  Missing evidence is never converted into a zero-duration option.

Memory evidence is obtained in this order: an enabled hardware-matched ML
predictor; an exact batch observation; linear interpolation between the closest
lower and upper observations; then, only at the originally requested batch, a
solo profile or submitted estimate.  Runtime evidence comes from a matching
runtime profile, with an exclusive-backend profile as fallback, then a batch
observation containing seconds per epoch or step timing.  These source labels
and confidence values are persisted for auditing.

Option $b'$ dominates $b$ when

\begin{equation}
  \widehat m_{i,h}(b')\le\widehat m_{i,h}(b),\qquad
  \widehat p_{i,h}(b')\le\widehat p_{i,h}(b),
  \label{eq:dominance-weak}
\end{equation}

and at least one inequality is strict.  Dominated options are removed.  Let
$\mathcal P_{i,h}\subseteq\mathcal B_i$ denote the retained Pareto frontier.

For active job $j\in A(t)$, the current resolved batch is fixed; it is not
re-optimized.  The estimator obtains its remaining runtime at that batch, or
falls back to a job-level remaining-runtime profile.  Failure to obtain a
positive active-job runtime makes the entire candidate infeasible.

\subsection{Step 8: enforce compatibility and predicted memory}

The candidate is evaluated as the combined set

\begin{equation}
  H=A(t)\cup G.
\end{equation}

When $|H|>1$, every member must be explicitly packing-eligible, have a stable
packing signature, and allow backend $h$.  For every unordered pair
$\{i,j\}\subseteq H$, a stored pair profile rejects the group if that profile is
on cooldown or marked incompatible:

\begin{equation}
  \mathrm{Compatible}(H,h)
  =\bigwedge_{\{i,j\}\subseteq H}
   \left[
     \neg\mathrm{HasPairProfile}_{ij,h}
     \;\lor\;
     \bigl(\neg\mathrm{Cooldown}_{ij,h}\land
            \mathrm{PairCompatible}_{ij,h}\bigr)
   \right].
  \label{eq:pair-compatibility}
\end{equation}

An absent pair profile does not reject a time-aware group.  Unlike older
packing modes, this compatibility predicate intentionally does not use SM
utilization or a slowdown threshold.

The predicted memory budget is

\begin{equation}
  B_{\mathrm{pred}}
  =\rho_{\mathrm{pred}}M_{\mathrm{total}},
  \label{eq:predicted-budget}
\end{equation}

where $M_{\mathrm{total}}$ is configured device VRAM when present, otherwise
detected hardware VRAM, and $\rho_{\mathrm{pred}}=0.85$ in the current example.
For a batch-option vector
$\boldsymbol b=(b_i)_{i\in G}\in\prod_{i\in G}\mathcal P_{i,h}$, feasibility
requires

\begin{equation}
  M_{\mathrm{active}}(t)
  +\sum_{i\in G}\widehat m_{i,h}(b_i)
  \le B_{\mathrm{pred}}+10^{-9}.
  \label{eq:memory-feasibility}
\end{equation}

The $10^{-9}$ term is only a floating-point boundary tolerance.  The memory
quantity is predicted average used VRAM, not peak allocation.

\subsection{Step 9: calculate priority weights}

Let the minimum effective priority in the window be

\begin{equation}
  \pi_{\min}(t)=\min_{i\in W(t)}\pi_i(t).
\end{equation}

Each waiting-window job receives the non-negative relative weight

\begin{equation}
  w_i(t)=1+\eta\bigl(\pi_i(t)-\pi_{\min}(t)\bigr),
  \qquad i\in W(t).
  \label{eq:priority-weight}
\end{equation}

Subtracting $\pi_{\min}$ makes the least urgent window job have weight one and
preserves priority differences.  Active jobs normally do not belong to the
runnable window; the implementation assigns such a job the default weight
$\bar w_j=1$.  In general define

\begin{equation}
  \bar w_j=
  \begin{cases}
    w_j(t), & j\in W(t),\\
    1, & j\notin W(t).
  \end{cases}
  \label{eq:active-weight}
\end{equation}

\subsection{Step 10: construct the exclusive-anchor baseline}

Let $a$ be the mandatory anchor when one exists, otherwise the first job in the
window.  Among the anchor's feasible exclusive batch options, choose the
smallest predicted remaining runtime:

\begin{equation}
  p_a^{\mathrm{ex}}
  =\min_{b\in\mathcal P_{a,\mathrm{exclusive}}:\,
          \widehat m_{a,\mathrm{exclusive}}(b)\le B_{\mathrm{pred}}}
    \widehat p_{a,\mathrm{exclusive}}(b).
  \label{eq:exclusive-anchor-duration}
\end{equation}

If there are no batch-indexed options, a positive job-level exclusive runtime
estimate is used.  The normalization baseline is the incremental flow cost of
letting the exclusive anchor occupy the next drain interval:

\begin{equation}
  L_{\mathrm{ex}}
  =p_a^{\mathrm{ex}}\sum_{i\in W(t)}w_i(t).
  \label{eq:exclusive-baseline}
\end{equation}

If this baseline cannot be estimated, no time-aware comparison is possible.  On
an idle GPU the planner falls back to the anchor alone on \code{exclusive}; if
work is already active, it returns no dispatch.

\subsection{Step 11: evaluate rolling-horizon predicted flow cost}

For a feasible triple $(G,h,\boldsymbol b)$, let

\begin{equation}
  H=A(t)\cup G
\end{equation}

and assign completion offsets from the present decision time:

\begin{equation}
  d_j=
  \begin{cases}
    \widehat p_{j,h}(b_j), & j\in G,\\
    \widehat p^{\mathrm{remaining}}_{j,h}, & j\in A(t).
  \end{cases}
  \label{eq:completion-offsets}
\end{equation}

The predicted drain time is

\begin{equation}
  D(G,h,\boldsymbol b)=\max_{j\in H}d_j.
  \label{eq:drain-time}
\end{equation}

Selected and active jobs accrue their individual completion offsets.  Every
window job not selected is assumed to wait throughout the complete drain
interval.  Therefore the implemented one-horizon weighted flow cost is

\begin{equation}
  L_F(G,h,\boldsymbol b)
  =\sum_{j\in H}\bar w_j d_j
   +D(G,h,\boldsymbol b)
      \sum_{i\in W(t)\setminus H}w_i(t).
  \label{eq:flow-cost}
\end{equation}

The score is

\begin{equation}
  J(G,h,\boldsymbol b)
  =\frac{L_F(G,h,\boldsymbol b)}{\max(10^{-9},L_{\mathrm{ex}})}.
  \label{eq:normalized-score}
\end{equation}

The denominator is constant for all candidates in one call, so normalization
does not change their ordering.  It gives the logged value an interpretable
reference: $J<1$ predicts lower one-horizon weighted flow than running the
anchor exclusively, $J=1$ is equal, and $J>1$ is worse.  There is no separate
minimum-gain threshold in the implementation.

This model uses solo/batch-indexed remaining times directly.  It does not inflate
$d_j$ by a pairwise slowdown prediction.  Stored slowdown observations remain
passive evidence for reports and do not enter Equations~\eqref{eq:drain-time}--
\eqref{eq:normalized-score}.

\subsection{Step 12: bounded search over batch vectors}

For a group with at most $N_{\mathrm{exact}}$ new jobs, every vector in

\begin{equation}
  \mathcal V_{G,h}=\prod_{i\in G}\mathcal P_{i,h}
\end{equation}

is evaluated.  For larger groups, the batch-option search expands one job at a
time and retains only $B_{\mathrm{beam}}$ partial vectors using

\begin{equation}
  k^{\mathrm{option\mbox{-}beam}}(\boldsymbol b)
  =\left(
     \sum_i\widehat p_{i,h}(b_i),
     \sum_i\widehat m_{i,h}(b_i),
     (b_i)_i
   \right).
  \label{eq:option-beam}
\end{equation}

This beam key favors low summed runtime, then low memory, then the
lexicographically smaller batch vector.  The objective is evaluated only after
the retained full vectors have been constructed.

\subsection{Step 13: deterministic final selection}

Each feasible candidate is compared with the lexicographic key

\begin{equation}
\begin{split}
  K(G,h,\boldsymbol b)=\big(&J,
    \min_{i\in G}s_i,
    -|G|,
    \sum_{i\in G}\widehat m_{i,h}(b_i),\\
    &\operatorname{sort}(\mathrm{id}_i:i\in G),
    \operatorname{sort}((\mathrm{id}_i,b_i):i\in G),
    \operatorname{rank}(h),h\big).
  \label{eq:tie-key}
\end{split}
\end{equation}

The selected candidate is

\begin{equation}
  (G^*,h^*,\boldsymbol b^*)
  =\argminlex_{(G,h,\boldsymbol b)\ \mathrm{feasible}}
     K(G,h,\boldsymbol b).
  \label{eq:final-choice}
\end{equation}

Thus exact score ties prefer, in order: the group with the oldest member; more
new jobs; lower predicted candidate memory; lexicographically smaller job IDs;
lexicographically smaller batch assignments; and only then the earlier backend
in \code{backend\_priority}.  Backend priority is not a primary optimization
term.

The plan mode is \code{exclusive} for one idle job,
\code{concurrent\_group} for one job added beside active work,
\code{packed\_pair} for two new jobs, and \code{packed\_group} for three or more
new jobs.

\subsection{Step 14: no-feasible-candidate and runtime fallbacks}

If no feasible candidate exists:

\begin{itemize}[leftmargin=*]
  \item with active work, return no plan and wait for a later scheduler cycle;
  \item on an idle GPU, dispatch the anchor alone on \code{exclusive}; and
  \item if starvation supplied $a^*$, preserve it as the mandatory anchor in the
        audit record.
\end{itemize}

At launch time the supervisor again rejects an unknown or unavailable backend,
an exclusive overlap, or an overlap involving a disallowed active backend.  If
a non-exclusive backend launch raises an exception while there are no other
active runs, the service makes one operational fallback attempt: the first
selected job alone on \code{exclusive}.  This exception path is separate from
the mathematical ranking.

The planner also stores a deterministic fallback ordering of group members by

\begin{equation}
  \bigl(p_i,-\widehat p^{\mathrm{job}}_{i,h},
        -\widehat m_{i,h}(b_i),q_i\bigr),
  \label{eq:fallback-order}
\end{equation}

ascending.  Consequently, lower base priority sorts first; ties prefer longer
job-level predicted remaining runtime (or zero when absent), then higher
predicted memory, then FIFO sequence.  In the
current code this order is persisted with group state and profiles; it is not a
term in the time-aware objective.

\section{Live-memory admission hysteresis}
\label{sec:live-gate}

Predicted feasibility in Equation~\eqref{eq:memory-feasibility} and live
admission are distinct protections.  Let telemetry sample $k$ report used and
total memory $(U_k,M_k)$, with fraction

\begin{equation}
  x_k=\frac{U_k}{M_k}.
\end{equation}

For the samples in the trailing window of length $T_{\mathrm{live}}$, define

\begin{equation}
  \overline x(t)=\frac{1}{n(t)}\sum_{k\in\mathcal I(t)}x_k,
  \qquad
  \mathcal I(t)=\{k:t-T_{\mathrm{live}}\le t_k\le t\}.
  \label{eq:rolling-memory}
\end{equation}

With the example thresholds
$\rho_{\mathrm{stop}}=0.90$,
$\rho_{\mathrm{resume}}=0.85$, and
$T_{\mathrm{live}}=10$ seconds:

\begin{align}
  \mathrm{open}\to\mathrm{closed}
  &\quad\text{when a complete window exists and}\quad
  \overline x(t)\ge\rho_{\mathrm{stop}}, \\
  \mathrm{closed}\to\mathrm{open}
  &\quad\text{after }\overline x(t)\le\rho_{\mathrm{resume}}
  \text{ continuously for }T_{\mathrm{live}}.
\end{align}

If the average rises above the resume threshold during the second condition,
the below-threshold timer resets.  This hysteresis prevents rapid open/closed
oscillation.  The gate state and samples are persisted across scheduler restart.
It blocks only additional admission beside active work and never pauses or kills
active jobs.

\section{Worked scoring example}

Suppose the window contains jobs $1,2,3$ with effective priorities
$(3,1,1)$ and $\eta=0.10$.  Then

\begin{equation}
  (w_1,w_2,w_3)=(1.2,1.0,1.0).
\end{equation}

Job $1$ is the anchor and its fastest feasible exclusive remaining time is
$50$ seconds.  Therefore

\begin{equation}
  L_{\mathrm{ex}}=50(1.2+1+1)=160.
\end{equation}

Consider a feasible two-job pack $G=\{1,2\}$ with predicted offsets
$(d_1,d_2)=(40,45)$ seconds.  Its drain time is $D=45$, while job $3$ remains
unselected.  Equation~\eqref{eq:flow-cost} gives

\begin{equation}
  L_F=1.2(40)+1(45)+45(1)=138,
  \qquad J=\frac{138}{160}=0.8625.
\end{equation}

If a different batch vector predicts $(d_1,d_2)=(40,60)$, then

\begin{equation}
  L_F=1.2(40)+1(60)+60(1)=168,
  \qquad J=1.05.
\end{equation}

Provided both vectors satisfy backend, compatibility, and memory constraints,
the first vector wins because $0.8625<1.05$.  This example also shows why the
drain term matters: a longer co-run delays every unselected window job.

\section{Compact algorithm listing}

The following pseudocode summarizes the exact decision structure without the
logging and persistence details.

\begin{enumerate}[label=\textbf{\arabic*.},leftmargin=*]
  \item Obtain runnable $Q$, active $A$, backend availability, predicted active
        VRAM, live-gate state, and exclusive-probe reservation.
  \item If a probe is reserved, return no plan until $A=\varnothing$; then run
        the probe exclusively.
  \item Build $W$ from top-priority, oldest, starving, and probe sets.  Set the
        oldest starving job as mandatory anchor, otherwise use $W[0]$.
  \item If active work exists and live admission is closed or a drain is
        requested, return no plan.  Enforce the remaining parallel-job cap.
  \item Compute $w_i$ and $L_{\mathrm{ex}}$.  If the baseline is missing, use
        the exclusive idle fallback or return no plan when active.
  \item Generate bounded candidate subsets $G$; enforce the mandatory anchor.
  \item For every legal backend $h$, reject the combined active/new set if
        packing eligibility or pair history is incompatible.
  \item Generate five requested-batch-relative options, estimate memory and
        remaining runtime, and Pareto-prune each job's options.
  \item Enumerate or beam-search option vectors.  Reject vectors exceeding
        $B_{\mathrm{pred}}$ after active memory is included.
  \item Compute completion offsets, drain time, $L_F$, normalized score $J$, and
        audit metadata.
  \item Select the lexicographically minimum key in
        Equation~\eqref{eq:tie-key}.  If none exists, wait when active or run the
        anchor exclusively when idle.
  \item Revalidate the plan in the supervisor, launch it, persist selected
        batches and the objective breakdown, and repeat admission while legal.
\end{enumerate}

\section{Other configured scheduler modes}

The planner retains five alternatives to \code{parallel\_time\_aware}.  They
share the priority-FIFO queue and exclusive fallback, but use different group
search and scoring rules.

\subsection{Serial modes}

Both \code{serial\_basic} and \code{serial\_batch\_optimized} select the first
priority-FIFO job and return an exclusive plan.  Let
$\pi_i^{\mathrm{policy}}(t)$ denote the generic policy's effective priority,
whose aging clock uses \code{waiting\_since()} (and can therefore reflect a
resume boundary).  Within the placement planner, their dispatch decision is
identical:

\begin{equation}
  i^*=\arg\min_i\bigl(-\pi_i^{\mathrm{policy}}(t),q_i\bigr),
  \qquad G^*=\{i^*\},\quad h^*=\mathrm{exclusive}.
\end{equation}

\subsection{Parallel default: fixed-batch VRAM fill}

After compatibility and memory checks, the fixed-batch score is maximized:

\begin{equation}
  J_{\mathrm{fixed}}(G,h)
  =\frac{\sum_{i\in G}\widehat m_{i,h}(b_i^{\mathrm{resolved}})}{B_{\mathrm{legacy}}}
   +0.001|G|-0.0001\left(\frac{1}{|G|}\sum_{i\in G}\max(1,q_i)\right).
  \label{eq:fixed-score}
\end{equation}

Unlike the time-aware score, larger is better.  Candidate groups come from a
fixed prefix of the priority queue and are bounded by
\code{max\_packed\_jobs\_per\_gpu}, with optional three-way packing.

\subsection{Parallel batch optimized}

This mode searches candidate batches and maximizes memory fill under a target
budget:

\begin{equation}
  J_{\mathrm{batch}}(G,h,\boldsymbol b)
  =\frac{\sum_{i\in G}\widehat m_{i,h}(b_i)}
         {\rho_{\mathrm{target}}B_{\mathrm{legacy}}},
  \qquad
  \sum_{i\in G}\widehat m_{i,h}(b_i)
  \le\rho_{\mathrm{target}}B_{\mathrm{legacy}}.
  \label{eq:batch-optimized-score}
\end{equation}

It uses a Cartesian product for groups of at most three jobs.  For larger
groups it evaluates a baseline vector and vectors that maximize one member's
batch at a time.

\subsection{Parallel auto-pack}

Auto-pack can score singleton additions and groups.  It requires runtime-ready
multi-job members, applies a runtime-skew penalty, and fills toward either a
VRAM or SM target.  Its maximized score is

\begin{equation}
  J_{\mathrm{auto}}(G,h)
  =1-\mathrm{gap}(G,h)+0.01|G|-P_{\mathrm{runtime}}(G,h)
   +0.02\,\mathbf{1}[h\ne\mathrm{exclusive}].
  \label{eq:auto-score}
\end{equation}

For a VRAM target,

\begin{equation}
  \mathrm{gap}
  =\frac{|\rho_VM_{\mathrm{safe}}-(M_{\mathrm{active}}+M_G)|}
         {\max(1,\rho_VM_{\mathrm{safe}})},
\end{equation}

and candidates above the target are rejected.  The SM form replaces memory by
summed predicted SM utilization and uses an unnormalized absolute gap.  This is
the only current placement mode whose scoring consumes the runtime-skew
guardrail.

\section{Interpretation, guarantees, and limitations}

\begin{itemize}[leftmargin=*]
  \item \textbf{Starvation protection is hard.}  Once the timeout is crossed,
        the oldest starving job must appear in every candidate and is the idle
        exclusive fallback.
  \item \textbf{Priority is relative.}  Priority changes window membership and
        the weighted-flow cost; it does not require strict dispatch order before
        the starvation threshold.
  \item \textbf{The horizon is local.}  Unselected jobs are charged one drain
        interval, but the scheduler does not recursively optimize all later
        dispatches.  It replans after state changes.
  \item \textbf{Search is bounded.}  Exact enumeration is used only at the
        configured small-job limit.  The two deterministic beams bound planning
        latency but can omit the global optimum.
  \item \textbf{Memory safety is predictive plus reactive.}  Candidate average
        VRAM must fit the predicted budget, while sustained live memory can close
        further overlap admission.  Neither is a proof against instantaneous
        peaks or prediction error.
  \item \textbf{Slowdown is not optimized.}  Co-run slowdown observations may be
        recorded after execution, but they do not affect the current flow-only
        decision unless they have already marked a pair incompatible or on
        cooldown.
  \item \textbf{Time-aware work is non-preemptive.}  Urgent jobs and probes wait
        for drain boundaries.  This preserves active training progress but can
        increase urgent-job latency.
  \item \textbf{Early stopping is orthogonal.}  It may shorten realized runtime
        at epoch-safe points, but the placement objective uses the remaining
        runtime estimates available at the decision instant.
\end{itemize}

\section{Current example parameters}

The repository's \code{configs/scheduler.example.yaml} currently supplies the
following time-aware values.  The algorithm remains parameterized, and deployed
configuration or detected hardware may differ.

\begin{center}
\begin{tabular}{@{}ll@{}}
\toprule
Parameter & Example value \\
\midrule
Mode & \code{parallel\_time\_aware} \\
Backend priority & \code{mps}, \code{stream}, \code{cuda\_process}, \code{exclusive} \\
Parallel job cap & unlimited (\code{null}) \\
Priority-window size $K$ & $8$ \\
Oldest-window size $L$ & $4$ \\
Starvation timeout & $1800$ s \\
Exact-search job limit & $3$ \\
Beam width & $64$ \\
Predicted-memory fraction & $0.85$ \\
Live stop/resume fractions & $0.90/0.85$ \\
Live averaging window & $10$ s \\
Priority weight $\eta$ & $0.10$ \\
Batch exponent offsets & $[-2,-1,0,1,2]$ \\
Objective version & \code{time\_v3\_flow\_only} \\
Exclusive probe policy & enabled, drain without preemption \\
\bottomrule
\end{tabular}
\end{center}

\section{Code-to-equation traceability}

\begin{longtable}{p{0.40\textwidth}p{0.52\textwidth}}
\toprule
Implementation location & Formalized behavior \\
\midrule
\endhead
\code{placement\_planner.py:74--118} & Effective priority, union planning window, and mandatory starvation anchor; Equations~\eqref{eq:effective-priority}--\eqref{eq:mandatory-anchor}. \\
\code{placement\_planner.py:120--142} & Exclusive-anchor normalization baseline; Equations~\eqref{eq:exclusive-anchor-duration}--\eqref{eq:exclusive-baseline}. \\
\code{placement\_planner.py:144--258} & Probe handling, active-work gates, backend enumeration, selection, fallbacks, and priority weights. \\
\code{candidate\_generator.py:63--107} & Immutable requested-batch proposals and clipping; Equations~\eqref{eq:batch-proposals}--\eqref{eq:batch-clipping}. \\
\code{candidate\_generator.py:156--191} & Exact subset generation and group beam; Equations~\eqref{eq:exact-groups}--\eqref{eq:group-beam}. \\
\code{resource\_estimator.py:43--54} & Time-aware predicted-memory budget; Equation~\eqref{eq:predicted-budget}. \\
\code{resource\_estimator.py:81--125} & Batch option bundle and remaining runtime; Equation~\eqref{eq:remaining-runtime}. \\
\code{resource\_estimator.py:127--288} & Evidence fallback order and Pareto dominance; Equation~\eqref{eq:dominance-weak}. \\
\code{compatibility.py:95--112} & Time-aware eligibility and pair-history rejection; Equation~\eqref{eq:pair-compatibility}. \\
\code{time\_objective.py:30--155} & Combined active/new feasibility, flow cost, normalized score, and candidate construction; Equations~\eqref{eq:memory-feasibility} and \eqref{eq:flow-cost}--\eqref{eq:normalized-score}. \\
\code{time\_objective.py:157--188} & Option-vector beam and final lexicographic tie key; Equations~\eqref{eq:option-beam}--\eqref{eq:final-choice}. \\
\code{telemetry.py:45--86} & Rolling live-memory hysteresis; Equation~\eqref{eq:rolling-memory}. \\
\code{service.py:1112--1173} & Exclusive reservation and repeated concurrent admission loop. \\
\code{supervisor.py:131--234} & Final backend availability and overlap validation. \\
\bottomrule
\end{longtable}

\end{document}
